\chapter{РЕЗУЛЬТАТЫ ЭКСПЕРИМЕНТА}

\section{Реализация доработки со структурированной памятью}

В рамках практической части работы была реализована доработка ClawdLab, связанная со структурированной памятью исследования. Архитектурно память реализована как отдельный слой поверх существующих сущностей ClawdLab. 
В реализованной схеме используются три уровня памяти. Рабочая память хранит текущий контекст выполнения: активную задачу, план, свежие результаты провайдеров и временные заметки. Эпизодическая память фиксирует завершенные фрагменты работы: краткие итоги попыток, обсуждений, критики и локальные уроки, полученные в ходе выполнения задачи. Каноническая память содержит устойчивые выводы лаборатории: принятые результаты, финальные заключения, правила лаборатории и обобщения, которые могут использоваться в следующих задачах. Такое разделение позволяет постепенно переводить информацию от сырого рабочего контекста к более устойчивому знанию.

Также была заложена логика продвижения памяти между уровнями. Для этого используется журнал операций продвижения, который фиксирует, из какого уровня в какой выполнялась дистилляция, какие записи были использованы на входе, какие записи были созданы на выходе и чем завершилась операция. Благодаря этому процесс обновления памяти становится наблюдаемым и воспроизводимым: можно понять, почему конкретный вывод оказался в каноническом слое и на какие данные он опирается.

\section{Итоговое сравнение метрик между конфигурациями}
